\documentclass{article}
\usepackage[portuguese]{babel}
\usepackage{listings}
\usepackage{xcolor}
\usepackage[margin=1.5cm]{geometry}

\definecolor{codegreen}{rgb}{0,0.6,0}
\definecolor{codegray}{rgb}{0.5,0.5,0.5}
\definecolor{codepurple}{rgb}{0.8,0,0.2}
\definecolor{backcolour}{rgb}{.8,.8,1}

\lstdefinestyle{mystyle}{
    backgroundcolor=\color{backcolour},   
    commentstyle=\color{codegreen},
    keywordstyle=\color{blue},
    numberstyle=\tiny\color{codegray},
    stringstyle=\color{codepurple},
    basicstyle=\ttfamily\footnotesize,
    breakatwhitespace=false,         
    breaklines=true,                 
    captionpos=b,                    
    keepspaces=true,                 
    numbers=left,                    
    numbersep=5pt,                  
    showspaces=false,                
    showstringspaces=false,
    showtabs=false,                  
    tabsize=2
}

\lstset{style=mystyle}
\begin{document}

\section{AA8}
Escreva um programa, em linguagem Python, para as seguintes tarefas:
\begin{enumerate}
    \item ler a dimensão n de uma matriz no console;
    \item registrar as entradas da matriz no console;
    \item imprimir a matriz original;
    \item imprimir a matriz na forma transposta;
    \item imprimir a matriz resultante da soma da matriz original com sua transposta;
    \item imprimir a matriz resultante do produto da matriz original com sua transposta;
    \item imprimir o valor do determinante da matriz original - aqui pode usar uma função pré-compilada;
    
\end{enumerate}
\subsection{}

Iniciamos o programa solicitando a ordem(dimensao) da nossa matriz ao usuario. Os elementos entao sao inseridos manualmente. 
\begin{lstlisting}[language=python]
ordem = int(input("Digite a dimensao da matriz A\n"))
matriz = []
for u in range(ordem):
    linha = []
    for v in range(ordem):
        linha.append(int(input(f"Digte o valor de A{u+1}{v+1}\n")))
    matriz.append(linha)
print("Matriz inicial A:")
printmatriz(matriz)
\end{lstlisting}
A matriz em nosso codigo e representada por uma lista de listas, onde cada lista interna representa uma linha da matriz.
Para imprimir nossa matriz no terminal definimos a funcao printmatriz() no inicio do nosso codigo, fazemos isso para poder reusar esse codigo.
\begin{lstlisting}[language=python]
def printmatriz(matriz):
    for linha in matriz:
        print(linha)
\end{lstlisting}

\subsection{Matriz transposta}
Para gerar nossa matriz transposta, fazemos as linhas se tornarem colunas
\begin{lstlisting}[language=Python]
transposta = []
for u in range(ordem):
    linha = []
    for v in range(ordem):
        linha.append(matriz[v][u])
    transposta.append(linha)
print("AT: Transposta de A:")
printmatriz(transposta)
\end{lstlisting}
Para cada elemento fazemos:
$A^T_{[v,u]} = A_{[u,v]}$

\subsection{Soma de matrizes}
Para somar duas matrizes apenas somamos os elementos de mesma posicao. $(A+B)_{[u,v]}=A_{[u,v]}+ B_{[u,v]}$

\begin{lstlisting}[language=Python]
soma = []
for u in range(ordem):
    linha = []
    for v in range(ordem):
        linha.append(matriz[u][v] + transposta[u][v])
    soma.append(linha)
print("Soma de A com AT")
printmatriz(soma)
\end{lstlisting}

\subsection{Produto de matrizes}
Para multiplicacao, cada elemento da matriz produto e o somatorio do produto de cada elemento da linha pela coluna correspondente.$P = AB$

\[
P_{[u,v]}= \sum_i A_{[u,i]} \cdot B_{[i,v]}
\]

\begin{lstlisting}[language=Python]
produto = []
for u in range(ordem):
    linha = []
    for v in range(ordem):
        elemento = 0
        for passo in range(ordem):
            elemento += matriz[u][passo] * transposta[passo][v]
        linha.append(elemento)
    produto.append(linha)
print("Produto de A com AT:")
printmatriz(produto)
\end{lstlisting}

\subsection{Determinantes}
Aqui simplesmentes usamos a biblioteca numpy.

\begin{lstlisting}[language=Python]
import numpy as np
#[...
#...]
determinante = np.linalg.det(matriz)
print(f"Determinante:{determinante}")
\end{lstlisting}

\subsection{Codigo completo:}
\begin{lstlisting}[language=Python]
import numpy as np

print("AA-08-Marcus Sousa\n\n")
deveParar = ""

def printmatriz(matriz):
    for linha in matriz:
        print(linha)

while (deveParar == ""):
    ordem = int(input("Digite a dimensao da matriz A\n"))
    matriz = []
    for u in range(ordem):
        linha = []
        for v in range(ordem):
            linha.append(int(input(f"Digite o valor na posisao {u+1},{v+1} \n")))
        matriz.append(linha)
    print("Matriz inicial A:")
    printmatriz(matriz)

    transposta = []

    for u in range(ordem):
        linha = []
        for v in range(ordem):
            linha.append(matriz[v][u])
        transposta.append(linha)
    print("AT: Transposta de A:")
    printmatriz(transposta)

    soma = []
    for u in range(ordem):
        linha = []
        for v in range(ordem):
            linha.append(matriz[u][v] + transposta[u][v])
        soma.append(linha)
    print("Soma de A com AT")
    printmatriz(soma)

    produto = []
    for u in range(ordem):
        linha = []
        for v in range(ordem):
            elemento = 0
            for passo in range(ordem):
                elemento += matriz[u][passo] * transposta[passo][v]
            linha.append(elemento)
        produto.append(linha)
    print("Produto de A com AT:")
    printmatriz(produto)

    determinante = np.linalg.det(matriz)
    print(f"Determinante:{determinante}")

    deveParar = input("Deseja continuar?(enter para continuar, qualquer tecla para parar)\n")

print("Ate mais")
\end{lstlisting}

\end{document}
